\documentclass[11pt]{article}

\usepackage[margin=1in]{geometry}
\usepackage{amsmath,amssymb}
\usepackage{booktabs}
\usepackage{longtable}
\usepackage{array}
\usepackage{tabularx}
\usepackage{xcolor}
\usepackage{xurl}
\usepackage{hyperref}
\usepackage{enumitem}
\usepackage{microtype}
\usepackage{fancyhdr}

\hypersetup{
  colorlinks=true,
  linkcolor=blue!55!black,
  urlcolor=blue!65!black,
  citecolor=green!40!black,
  pdftitle={Dual-Fuel Reactor Network Technical Reference},
  pdfauthor={fuelnozzle project}
}

\pagestyle{fancy}
\fancyhf{}
\lhead{Dual-Fuel Reactor Network Reference}
\rhead{Version 0.1.0}
\cfoot{\thepage}

\newcommand{\code}[1]{\nolinkurl{#1}}
\newcommand{\srcline}[2]{\nolinkurl{#1:#2}}
\newcommand{\unit}[1]{\,\mathrm{#1}}
\setlength{\emergencystretch}{3em}
\newcommand{\warningbox}[1]{%
  \begin{center}
  \fcolorbox{orange!60!black}{orange!8}{%
    \parbox{0.91\textwidth}{\textbf{Engineering caution.} #1}}
  \end{center}}

\title{Dual-Fuel Reactor Network Technical Reference\\
\large Jet-A and LNG Combustor Modeling for Nozzle and Architecture Design\\
\large Beginner-Oriented Equations, Algorithms, Source Locations, and Limitations}
\author{fuelnozzle project}
\date{Draft, in progress}

\begin{document}
\maketitle

\begin{abstract}
This report explains the equations and software implementation of the reactor-network
extension to the \code{fuelnozzle} package, for a reader with no prior background in
combustion, chemical kinetics, spray physics, or Python. The extension takes the output
of the Jet-A pressure-swirl and flashing-LNG nozzle models and carries it into the
combustor, so that emissions, temperature, mixing, and autoignition risk can be estimated
without a three-dimensional simulation. Its purpose is to support the design of a
dual-fuel combustor that burns Jet-A during landing and take-off and LNG during cruise,
with separate injector hardware for each fuel and a shared combustor. Every major
equation is connected to the Python function that implements it, and physical law,
numerical approximation, and empirical calibration are kept visibly separate throughout.
This is documentation for a reduced-order screening and design-exploration tool, not a
certification method.
\end{abstract}

\tableofcontents
\newpage

\section{Status of This Document}

This document is written alongside the implementation, section by section, as each
module is completed. Sections marked \emph{pending} are planned but not yet written,
and their absence means the corresponding code does not exist yet either.

The companion documents are:

\begin{itemize}
  \item \code{docs/CRN_PLAN.md}, the design intent, task list, and risk register;
  \item \code{docs/CRN_IMPLEMENTATION_LOG.md}, the build record, deviations from the
        plan, and verification evidence;
  \item \code{docs/technical_reference.tex}, the equivalent reference for the nozzle
        models that this extension consumes;
  \item \code{mech/README.md}, the kinetic mechanism registry and provenance.
\end{itemize}

\section{How to Read This Report}

Each section separates four different kinds of statement. A \emph{physical law} is a
conservation rule, such as mass entering a steady reactor equalling mass leaving it. A
\emph{numerical approximation} is how the computer solves that law, such as replacing a
long tube by several small stirred tanks. A \emph{calibration} is a number fitted to data.
An \emph{unmodeled effect} is deliberately absent. Keeping these labels separate prevents a
software test from being mistaken for proof that a model represents real hardware.

Source references use a file and line range, for example
\srcline{src/fuelnozzle/crn/autoignition.py}{56--63}. Function names are included so a
reader can search the source even when later edits move a line.

\section{Basic Vocabulary for a New Reader}

A \emph{reactor} is a simplified region in which the gas is represented by one average
temperature and composition. A \emph{reactor network} connects several such regions with
mass flows. Residence time means how long material remains in a region on average.
Equivalence ratio compares the actual fuel-to-air ratio with the amount needed for complete
combustion: one is stoichiometric, below one is lean, and above one is rich. Ignition delay
is the time between preparing a mixture and its rapid self-heating. An emission index is
grams of pollutant per kilogram of fuel.

A reactor network does not resolve three-dimensional jets, wakes, droplets, or wall films
the way computational fluid dynamics can. It is fast because those details are represented
by a few zones and calibrated connections. Its answers are only as physical as those zones,
flows, and calibrations.

\section{Why a Reactor Network, and What It Cannot Do}

The network is useful for screening many candidate designs before expensive CFD and rig
tests. It can conserve species, integrate finite-rate chemistry, and expose how residence
time or air staging changes a trend. It cannot discover a liner geometry, mixing pattern,
flame attachment, thermoacoustic mode, or spray-wall interaction without additional models
and evidence. The current design evaluator is therefore a prototype screening path, as
stated in \srcline{src/fuelnozzle/crn/evaluate.py}{1--14}.

\section{Plain-English Implementation Change Log}

\subsection{Stage 0: an unavailable safety calculation is not safe}

Earlier versions used the Python value \code{None} both when an ignition calculation did
not produce a delay and when the requested condition was outside the table. They then called
that result safe. A novice can think of this as a thermometer displaying nothing and the
software interpreting the blank display as ``not hot.'' A blank display does not establish
that.

\code{AutoignitionVerdict} now includes \code{UNKNOWN} at
\srcline{src/fuelnozzle/crn/autoignition.py}{56--63}. The function
\code{autoignition_margin} returns that verdict and an error when no delay is available at
\srcline{src/fuelnozzle/crn/autoignition.py}{302--323}. The optimizer-facing function
\code{evaluate_objectives} converts the unknown state into a constraint violation at
\srcline{src/fuelnozzle/crn/objectives.py}{87--119}. This does not claim the passage has
failed physically. It means the tool lacks enough evidence to accept it.

The more general implementation is \code{GateStatus}, \code{GateResult}, and
\code{aggregate_gate_status} in
\srcline{src/fuelnozzle/crn/status.py}{1--32}. A failed gate dominates an unknown gate, and
an unknown gate prevents a pass. \code{PointResult} and \code{DesignResult} expose separate
computational and acceptance states in
\srcline{src/fuelnozzle/crn/evaluate.py}{60--115}. Thus a numerical solver may finish
successfully while the design remains unknown because spray fidelity or validation evidence
is absent. The old \code{feasible} name remains as a compatibility property, but it is true
only for an acceptance pass.

The module-level explanation in \code{evaluate.py} and the printed example were also
changed to call the current all-vapour calculations diagnostics rather than validated
design answers. The test
\code{test_unavailable_ignition_delay_is_unknown_and_fail_closed} in
\srcline{tests/test_crn_autoignition_thermal.py}{268--290} guards this behavior. The
equation and evidence status for every design-driving model is tracked in
\code{docs/CRN_EQUATION_REGISTER.md}.

\subsection{Stage 1: one mission description and one shared liner}

Earlier CRN examples could repeat pressure, temperature, and flow values in a second
\code{MissionPoint}. Two copies of the same fact can quietly disagree. The function
\code{mission_point_from_operating} now converts the canonical \code{OperatingPoint} into a
CRN point and retains both the original point and its named pressure stations; see
\srcline{src/fuelnozzle/crn/mission.py}{15--42} and
\srcline{src/fuelnozzle/crn/design.py}{273--290}. The
\code{MissionProfile.from_icao_lto} check at
\srcline{src/fuelnozzle/crn/mission.py}{52--77} requires one Jet-A point for each of the four
ICAO modes---takeoff, climb-out, approach, and idle---and an explicit duration for every mode.
\code{MissionProfile.from_cruise} at
\srcline{src/fuelnozzle/crn/mission.py}{79--91} accepts one or more LNG cruise points.

\code{resolve_pressure_stations} at
\srcline{src/fuelnozzle/operating.py}{98--123} gives separate names to compressor discharge,
dome, combustor exit, and fuel-nozzle inlet pressure. The liner loss is the user-given
fraction of compressor-discharge pressure. The current first-order model takes dome pressure
equal to compressor discharge and subtracts the liner loss to obtain combustor-exit pressure.
It also checks that the active pump can reach chamber pressure plus the requested nozzle
drop. These are transparent bookkeeping assumptions, not a prediction of the detailed
pressure field.

A full annular combustor contains repeated fuel-injector ``cups.'' The
\code{SectorDefinition} conversions at
\srcline{src/fuelnozzle/crn/hardware.py}{23--41} multiply engine-total air and fuel flows by
the fraction of cups represented. Scaling both flows by the same number preserves their
ratio. This assumes identical cups and no circumferential distortion.

The new \code{SharedLinerGeometry} at
\srcline{src/fuelnozzle/crn/hardware.py}{110--167} contains one set of zone volumes, air-hole
effective areas, and a cooling-air destination used for both fuels.
\code{DualFuelHardware} at
\srcline{src/fuelnozzle/crn/hardware.py}{170--190} puts separate Jet-A and LNG passages inside
that liner. Their fixed effective areas determine how the dome air divides; the optimizer
cannot claim a different split without changing hardware. In
\code{DesignEvaluator.\_build}, \srcline{src/fuelnozzle/crn/evaluate.py}{184--231}, supplied
hardware replaces the legacy optimized air fractions and volumes. Prescribed fractions are
still allowed for comparison with data, but \code{DualFuelHardware.\_\_post\_init\_\_}
requires a calibration identifier so an unexplained number cannot masquerade as geometry.

\subsubsection{How air flow through a hole is estimated}

\code{AirAdmission.mass\_flow\_kg\_s} at
\srcline{src/fuelnozzle/crn/hardware.py}{58--93} uses the standard compressible ideal-gas
orifice equation. In plain language, more effective hole area or upstream pressure passes
more air, while hotter air is less dense and passes less mass. When the downstream-to-upstream
pressure ratio falls below the critical value, the throat reaches the speed of sound and
further lowering downstream pressure no longer raises this idealized mass flow; this is
called choking.

The equation assumes steady, one-dimensional, isentropic air with constant
$\gamma=1.4$ and $R=287.05$ J/(kg K). A discharge coefficient accounts approximately for
contraction and losses. \code{SharedLinerGeometry.area\_derived\_split} at
\srcline{src/fuelnozzle/crn/hardware.py}{136--167} evaluates all five admissions between the
same prescribed upstream and downstream pressures and normalizes their flows to fractions.
It does \emph{not} yet solve the pressure loss from the geometry, force the absolute sum to
equal engine air flow, predict axial pressure variation, or describe jets penetrating a
crossflow. Cold-flow or sector-rig data are therefore still required, and design acceptance
remains unknown.

\subsubsection{How a declared LNG mixture reaches chemistry}

Property software calls methane ``Methane,'' whereas the chemical mechanism calls it
\code{CH4}. \code{LNGComposition.cantera\_mole\_fractions} at
\srcline{src/fuelnozzle/models.py}{69--85} performs the explicit name conversion for methane,
ethane, propane, and nitrogen. \code{MechanismRegistry.with\_lng\_composition} at
\srcline{src/fuelnozzle/crn/chemistry.py}{195--208} then copies those mole fractions into both
LNG chemistry roles. \code{validate\_mechanism} at
\srcline{src/fuelnozzle/crn/chemistry.py}{223--273} checks the mapped names against the loaded
mechanism. An unsupported component stops with an error instead of being dropped, because
dropping part of a fuel changes its heating value, stoichiometric air requirement, flame
temperature, and emissions.

\subsection{Stage 2: making the reactor graph match the stated equations}

\subsubsection{A plug-flow label now creates plug flow numerically}

A plug-flow region represents gas moving downstream without mixing back upstream. Cantera's
simultaneous recirculating network cannot insert a spatial plug-flow object directly, so the
usual approximation is a row of small stirred tanks. Material leaving tank one enters tank
two, then tank three, and so on. More tanks reduce the artificial axial mixing.

Earlier, \code{ReactorKind.PFR} stored a segment count but the solver made one large stirred
reactor. \code{\_expand\_plug\_flow\_reactors} now constructs the chain at
\srcline{src/fuelnozzle/crn/network.py}{392--480}. It preserves the physical zone's total
volume and heat loss, sends normal inlets to the first segment, and takes the outlet from the
last. \code{NetworkSolution.zone} and \code{zone\_residence\_time\_s} at
\srcline{src/fuelnozzle/crn/network.py}{344--375} let a reader recover all segments and their
total residence time. The production path is refined at 2, 4, and 8 segments in
\srcline{tests/test_crn_verification.py}{207--235}. Segment convergence is a numerical check,
not evidence that a real combustor has plug flow.

\subsubsection{A flow can no longer become negative and then disappear}

At steady state, mass entering every reactor must equal mass leaving it. Earlier, a
least-squares correction could make a declared flow negative. The Cantera assembly then
skipped that negative controller, so the graph actually solved was not the graph reported as
balanced.

\code{minimum\_norm\_mass\_correction} at
\srcline{src/fuelnozzle/crn/network.py}{100--318} retains its old name for compatibility but
now solves a constrained problem. Every corrected flow must be nonnegative. A fixed measured
or prescribed recycle cannot change. When more than one correction can close the graph, each
adjustable flow needs an uncertainty; changing a poorly known flow costs less than changing a
well-known one. Boundary flow must balance both globally and inside every disconnected
component. If the directed arrows cannot carry the required mass, the function stops instead
of reversing an arrow. New code may use its clearer alias \code{close\_internal\_flows}.

\code{Architecture.fixed\_internal\_flows} identifies physical recycle edges at
\srcline{src/fuelnozzle/crn/templates.py}{71--87}. \code{DesignEvaluator.evaluate\_point} and
\code{solve\_coupled} pass those identities into network assembly; see
\srcline{src/fuelnozzle/crn/evaluate.py}{251--258} and
\srcline{src/fuelnozzle/crn/coupling.py}{67--130}. Fuel vapor is therefore present before
closure. Its mass increases the downstream through-flow without being spread arbitrarily
around the recycle loop.

\subsubsection{Moles, mass, volume, and residence time}

Mole fraction answers the question ``what fraction of the molecules are this species?''
Mass flow answers ``how many kilograms pass each second?'' They cannot be multiplied across
different gases and added directly, because a kilogram of hydrogen contains many more
molecules than a kilogram of carbon dioxide. \code{\_initial\_mixture} at
\srcline{src/fuelnozzle/crn/network.py}{854--883} first divides each stream's mass flow by its
molecular weight to obtain molar flow, then accumulates each species. It separately combines
stream enthalpy using mass flow, which is the correct weighting for energy.

Residence time is
\[
  \tau = \frac{m}{\dot m} = \frac{\rho V_{\mathrm{physical}}}{\dot m}.
\]
A Cantera constant-pressure zero-dimensional reactor naturally holds its initial mass and lets
its volume change as it heats. That is useful for a piston but wrong for a fixed combustor
zone. The time-marching loop in \code{CombustorNetwork.solve},
\srcline{src/fuelnozzle/crn/network.py}{590--835}, re-imposes the declared physical volume
and reinitializes Cantera after every steady iteration. \code{\_extract} at
\srcline{src/fuelnozzle/crn/network.py}{925--954} reports Cantera's actual mass and volume,
checks $m=\rho V$, and calculates residence time from actual inventory rather than from a
different bookkeeping value. This volume reset is a numerical steady-state constraint; it
must not be interpreted as a real transient piston motion.

\subsubsection{What ``converged'' now means}

A temperature that has stopped moving does not guarantee trace species or stored mass have
settled. \code{ConvergenceReport} at
\srcline{src/fuelnozzle/crn/network.py}{89--98} records the largest change in temperature,
every species mass fraction, reactor mass, enthalpy inventory, physical volume, and the scaled
state derivative. The acceptance checks are in
\srcline{src/fuelnozzle/crn/network.py}{700--770}. After marching, the independent
\code{\_element\_balance\_error} and \code{\_energy\_balance\_error} calculations at
\srcline{src/fuelnozzle/crn/network.py}{956--1024} require atoms and steady-flow enthalpy to
close. Chemical reactions rearrange atoms and release chemical energy; they do not create
atoms or total enthalpy.

Combustion equations can have both an unlit cold steady state and a burning hot steady state
at one boundary condition. \code{InitializationBranch} at
\srcline{src/fuelnozzle/crn/network.py}{68--74}, \code{solve\_branches} at
\srcline{src/fuelnozzle/crn/network.py}{836--852}, and \code{solve\_continuation} at
\srcline{src/fuelnozzle/crn/network.py}{1027--1048} preserve that distinction. Continuation
uses one solution to seed the next point on an identical expanded topology. A hot mathematical
branch does not by itself prove that a real flame anchors, and a cold branch does not by
itself locate lean blowout; those require flameholding physics and rig evidence.

\subsubsection{Cooling and prescribed heat}

\code{\_cooling\_inlet} at
\srcline{src/fuelnozzle/crn/templates.py}{188--211} now reads the declared destination.
Primary cooling rejoins the primary/mixing zone, dilution cooling enters the beginning of the
post-flame chain, and exit cooling joins its final segment. This matters because upstream air
can affect reaction and NO formation for longer than air added at the exit.
\code{ReactorSpec.validate\_kind\_consistency} at
\srcline{src/fuelnozzle/crn/reactors.py}{66--88} also rejects a positive
\code{heat\_loss\_w} without a calibration or physical-model identifier. That makes prescribed
heat traceable, but it does not make the number predictive. A wall/coolant model and liner
thermal data remain Stage 4 requirements.

\subsection{Stage 3: following real liquid fuel from the nozzle}

\subsubsection{Why the old all-vapour shortcut was not enough}

The earlier design loop placed every kilogram of fuel into the first reactor as gas. That is
quick, but a liquid droplet may spend much of the combustor heating and evaporating. Fuel still
inside a droplet cannot burn. Treating it as gas too early can therefore predict a flame that is
too hot, too stable, or in the wrong place.

\code{DesignEvaluator} now accepts a \code{spray_model} in
\srcline{src/fuelnozzle/crn/evaluate.py}{181--236}. The callback must return the actual
\code{SprayBoundary} and a liquid-property provider for each mission point.
\code{DesignEvaluator.evaluate_point}, at
\srcline{src/fuelnozzle/crn/evaluate.py}{296--696}, checks the fuel name and total nozzle flow,
then calls \code{solve_coupled}. If no callback is supplied, the old vapour-only calculation is
still available for fast diagnosis, but its acceptance status stays \code{UNKNOWN}.

For LNG, a pressure drop can separate the mixture into a vapour rich in light components and a
liquid rich in heavier components. \code{ThermodynamicState} and the nozzle CFD boundary now
carry both equilibrium phase compositions. \code{lng_spray_boundary}, at
\srcline{src/fuelnozzle/crn/spray_source.py}{230--347}, converts names such as
\code{Methane} into chemistry names such as \code{CH4}. The nozzle's finite-rate quality sets
how much mass is vapour, while the retained equilibrium compositions state the available
thermodynamic phase estimate. This distinction is reported because finite-rate fractionation
still needs data.

\subsubsection{Breakup now includes damping and is actually invoked}

The Taylor analogy imagines a droplet as a tiny damped spring. Air pushes it out of shape,
surface tension pulls it round, and viscosity damps its motion. The damped response used by
\code{taylor_analogy_breakup}, at
\srcline{src/fuelnozzle/crn/droplets.py}{139--238}, is
\[
 y(t)=y_f\left[1-e^{-\beta t}
 \left(\cos\omega_d t+\frac{\beta}{\omega_d}\sin\omega_d t\right)\right].
\]
Breakup occurs when $y=1$. A bounded root solve finds the first crossing. This corrects the
previous undamped timing. \code{\_march_droplets}, at
\srcline{src/fuelnozzle/crn/coupling.py}{383--508}, calls this equation only when the nozzle
boundary requests aerodynamic breakup. LNG flash breakup remains separate because it is driven
from inside the liquid, not by air loading.

\subsubsection{Mass, speed, vapour inhibition, and heat}

The droplet integrator previously advanced radius and inferred remaining mass from radius
cubed. That works only when liquid density is constant. Heating changes density, so the old
method could make mass appear or disappear. \code{integrate_droplet}, at
\srcline{src/fuelnozzle/crn/droplets.py}{437--581}, now advances mass directly and calculates
radius from
\[
 r=\left(\frac{3m}{4\pi\rho_l(T,P)}\right)^{1/3}.
\]

It also advances droplet velocity. The drag acceleration is
\[
 \frac{du_d}{dt}=-\frac{3C_D\rho_g}{8\rho_l r}
 (u_d-u_g)\left|u_d-u_g\right|,
\]
with the standard Schiller--Naumann drag coefficient. Local gas velocity comes from declared
spray-path length divided by gas residence time. The final speed of one zone becomes the
initial speed of the next.

Evaporation slows when fuel vapour already surrounds a droplet. \code{\_gas_state}, at
\srcline{src/fuelnozzle/crn/coupling.py}{511--538}, now sums the local fuel-species mass
fractions instead of always reporting zero. Finally, the droplet ODE integrates convective heat
itself. Multiplying one droplet's integrated heat by its number rate gives watts removed from
the gas. This avoids rebuilding a variable-property enthalpy change from only two endpoint
temperatures.

\code{CoupledSolution}, at
\srcline{src/fuelnozzle/crn/coupling.py}{62--73}, reports evaporated fuel, vapour injected by
the nozzle, liquid carryover, the residual of their mass balance, and integrated droplet heat.
Carryover is removed from gaseous exhaust flow before emission index is calculated. A closed
software ledger is verification; high-pressure spray and wall-impingement experiments are
still required for validation.

\subsection{Stage 4: heating LNG without confusing the feed line and nozzle}

LNG first travels through a pressurized feed line and only then drops through the nozzle.
Heating and pressure loss occur along the first path. Flashing occurs during the second,
approximately constant-enthalpy path. Using chamber pressure for the whole heating calculation
mixes these two physical events.

\code{feed_path_heat_budget}, at
\srcline{src/fuelnozzle/crn/thermal.py}{74--102}, checks the actual path produced by
\code{solve_lng_feed_line}: the heat rate must equal mass flow times the outlet-minus-inlet
enthalpy. Each feed-line station now retains quality and both phase compositions.

\code{heat_sink_budget}, at
\srcline{src/fuelnozzle/crn/thermal.py}{105--218}, uses the target pressure-temperature state,
the bubble enthalpy, and the dew enthalpy. A mixture boils over a temperature \emph{range}, not
at one temperature. A target inside that glide therefore receives a quality between zero and
one and only part of the total latent duty. \code{thermal_window}, at
\srcline{src/fuelnozzle/crn/thermal.py}{343--474}, evaluates heating at feed pressure, then
separately compares the target temperature with chamber-pressure saturation.

An engine cannot supply unlimited heat. \code{HeatSourceModel}, at
\srcline{src/fuelnozzle/crn/thermal.py}{231--274}, requires source temperature, source mass
flow, heat capacity, exchanger effectiveness, minimum pinch temperature, and fuel-side
pressure loss. Its available heat is the smaller of the effectiveness capacity and pinch
capacity. Without a hardware evidence identifier, this calculation remains \code{UNKNOWN}.

Inverse calculations must not hide impossible requests. \code{solve_temperature_for_duty}, at
\srcline{src/fuelnozzle/crn/thermal.py}{644--681}, and
\code{fuel_temperature_for_target_superheat} raise
\code{InfeasibleThermalTargetError} when the answer is outside the bracket or would boil
upstream. They no longer return the nearest limit as though it were the requested answer.

The idle passage matters too. \code{idle_circuit_screen}, at
\srcline{src/fuelnozzle/crn/thermal.py}{501--617}, still checks Jet-A coking temperature and LNG
vapour lock. It now separately requires soak duration, positive purge flow and time, and
demonstrated restart before its transient gate can pass. A steady wall temperature alone
cannot establish that a switched-off circuit will restart.

\subsection{Stage 5: what the chemistry and operability calculations actually prove}

\code{DesignEvaluator.evaluate_point} calls \code{validate_mechanism} before solving and stores
the mechanism file and provenance in \code{PointResult}; see
\srcline{src/fuelnozzle/crn/evaluate.py}{72--105}. A mechanism can load successfully yet still
be outside its declared pressure or temperature domain. Separate applicability and held-out
validation gates prevent ``the file ran'' from becoming ``the chemistry is accurate.''

\code{IgnitionDelayTable}, at
\srcline{src/fuelnozzle/crn/autoignition.py}{244--347}, now interpolates logarithmic ignition
delay in three dimensions: inverse temperature, pressure, and equivalence ratio. It never
extrapolates outside the table. \code{IgnitionEvidenceState}, at
\srcline{src/fuelnozzle/crn/autoignition.py}{67--72}, distinguishes:
\begin{itemize}
  \item an interpolated finite ignition delay;
  \item a censored lower bound, meaning no ignition happened before the integration stopped;
  \item unavailable evidence, meaning the requested state is outside the table.
\end{itemize}
A censored lower bound can establish a safe numerical margin only when even that lower bound
exceeds the required margin. \code{IgnitionMarker}, at
\srcline{src/fuelnozzle/crn/autoignition.py}{75--77}, permits both a fixed temperature rise and
the time of maximum temperature-rise rate so marker sensitivity can be measured.

\code{laminar_flame_speed}, at
\srcline{src/fuelnozzle/crn/autoignition.py}{583--607}, solves a freely propagating Cantera
flame where supported. \code{flashback_screen}, at
\srcline{src/fuelnozzle/crn/autoignition.py}{490--580}, compares passage velocity with an
estimated turbulent flame speed. A safe number cannot pass acceptance without a calibration
identifier for the actual passage.

\code{continuation_lean_limit_screen}, at
\srcline{src/fuelnozzle/crn/operability.py}{36--133}, follows the hot mathematical CRN branch
from rich to lean and brackets where that branch becomes extinguished.
\warningbox{This is a numerical extinction screen, not lean blowout. Real blowout also depends
on flame anchoring, turbulence, recirculation, heat loss, and disturbances. It remains
\code{UNKNOWN} until pressure-matched rig calibration exists.}
Transient ignition, relight, fuel switching, and thermoacoustics likewise remain explicit
external gates rather than invented reduced-order predictions.

\subsection{Stage 6: optimizing mission quantities instead of internal-reactor spreads}

Jet-A landing and take-off NOx is assembled as
\[
 D_p/F_{oo}=\frac{\sum_i EI_{\mathrm{NOx},i}\,\dot m_{f,i}\,t_i}
 {F_{oo}},
\]
where the four terms are takeoff, climb-out, approach, and idle. The corrected
\code{DesignResult.weighted_ei_nox} and \code{lto_emissions} are in
\srcline{src/fuelnozzle/crn/evaluate.py}{108--174}. Fuel flow times duration is fuel mass, so
long idle time no longer receives the same importance as a high-flow mode simply because it is
long. Rated thrust and all four modes are required before the value may drive optimization.

LNG results stay attached to each named cruise point through
\code{lng_cruise_ei_by_point}. Blank or duplicate cruise names are rejected by
\code{MissionProfile.from_cruise} at
\srcline{src/fuelnozzle/crn/mission.py}{79--93}.

The former ``mixing'' objective was merely the range of equivalence ratio across reactors in
series; that mostly measures combustion progress. The former ``exit spread'' was the
temperature range across every internal reactor; that is not an exit traverse. Both are now
unavailable (\code{NaN}) compatibility fields and are absent from the primary objective list.
A calibrated mixture-fraction distribution and mass-weighted parallel exit streams must be
supplied before those ideas return as objectives.

\code{evaluate_objectives}, at
\srcline{src/fuelnozzle/crn/objectives.py}{66--193}, retains each mode and cruise metric and
turns failed gates into constraints. \code{rank_key}, at
\srcline{src/fuelnozzle/crn/objectives.py}{196--210}, uses a declared lexicographic preference
only after feasibility; it no longer adds quantities with units of g/kN, g/kg, kelvin, and
dimensionless spread. The Pareto front remains the primary result.

\code{evaluate_uncertainty_ensemble}, at
\srcline{src/fuelnozzle/crn/uncertainty.py}{68--124}, evaluates explicit input, calibration,
mechanism, manufacturing, numerical, and model-form cases. It reports empirical objective
intervals and passes robust feasibility only when every category is represented and every case
passes. The software cannot supply the uncertainty distributions; those remain engineering
inputs backed by data.

\subsection{Stage 7: a report schema that says NO-GO when evidence is missing}

\code{EvidenceGrade} and \code{EvidenceRecord}, at
\srcline{src/fuelnozzle/crn/evidence.py}{14--49}, label each quantity verified, calibrated,
validated, extrapolated, unavailable, or missing. Verification shows that software solves its
stated equations. Validation shows, using independent data, that those equations predict
reality in a declared range. They are not interchangeable.

\code{build_conceptual_design_report}, at
\srcline{src/fuelnozzle/crn/evidence.py}{124--255}, gathers effective areas, pressure losses,
flows, volumes, residence times, thermal schedule, constraints, Pareto position, uncertainty,
and review evidence. Missing dimensions or injector ranges are written as \code{None}; a
reasonable-looking placeholder is never substituted.

A release is \code{GO} only when mission acceptance passes, all required predictive evidence is
validated, every uncertainty category is robustly feasible, and an independent technical
review is supplied. The present disposition is \textbf{NO-GO}: the repository does not contain
the engine cycle deck, installed geometry, nozzle/spray/thermal rig data, high-pressure
chemistry and emissions holdouts, exit traverse, transient operability evidence, uncertainty
inputs, or independent review needed for a dual-fuel conceptual-design recommendation.

\section{Kinetic Mechanisms}

\subsection{What a mechanism is}

A \emph{chemical kinetic mechanism} is a list of the molecules a fuel can turn into and
the reactions that convert one into another, each with a rate. Burning a hydrocarbon is
not one reaction but hundreds acting at once, so the mechanism is a file, not a formula.
The files used here contain 47 to 71 chemical species and 247 to 538 reactions.

Two properties of a mechanism matter for this work. First, whether it contains
\emph{nitrogen chemistry}: without species such as NO and NO$_2$, the mechanism cannot
produce oxides of nitrogen at all. Second, whether it contains \emph{low-temperature
chemistry}: the slow reactions that determine how long a fuel-air mixture can sit at
moderate temperature before it ignites on its own.

\subsection{Why Jet-A needs two mechanisms}

Neither available Jet-A file has both properties, and this tool needs both.

\begin{center}
\begin{tabular}{lccc}
\toprule
File & Species & Nitrogen chemistry & Low-temperature chemistry \\
\midrule
\code{A2NOx_skeletal}  & 71 & yes & no  \\
\code{A2NTCfast_ske}   & 47 & no  & yes \\
\code{gri30} (LNG)     & 53 & yes & n/a \\
\bottomrule
\end{tabular}
\end{center}

It would be convenient to use the nitrogen-bearing file for everything. The following
measurement, made during implementation, shows why that would be unsafe. The table gives
the ignition delay predicted by each Jet-A mechanism at conditions representative of a
premixing passage, expressed as the ratio between them.

\begin{center}
\begin{tabular}{lcccc}
\toprule
Pressure & 700 K & 800 K & 900 K & 1000 K \\
\midrule
20 atm & $510\times$ & $71\times$ & $9.3\times$ & $2.0\times$ \\
40 atm & $490\times$ & $86\times$ & $13.5\times$ & $2.9\times$ \\
\bottomrule
\end{tabular}
\end{center}

Compressor discharge temperature at landing and take-off is roughly 700 to 900~K.
In that range the high-temperature-only mechanism predicts an ignition delay one to two
orders of magnitude longer than the mechanism that includes low-temperature chemistry.

\warningbox{An ignition delay that is predicted too long makes a premixing passage look
safe when it is not. This tool therefore uses \code{A2NOx_skeletal} for the reactor
network and all NO$_x$ output, and \code{A2NTCfast_ske} only for ignition delay. Each
result records which mechanism produced it.}

\subsection{A distinction that is easy to get wrong}

Some hydrocarbons show a \emph{negative temperature coefficient}: a temperature band in
which raising the temperature makes ignition take \emph{longer}, contrary to intuition.
Measured for \code{A2NTCfast_ske}, that band is clear at 5~atm between roughly 800
and 950~K, weak at 10~atm, and gone by 20 to 40~atm.

It is tempting to conclude that low-temperature chemistry stops mattering at combustor
pressure. It does not. The $510\times$ measurements above were made at 20 and
40~atm, where the non-monotonic shape has already disappeared. What vanishes at
high pressure is the \emph{shape} of the curve, not the \emph{influence} of the
low-temperature reactions on the absolute delay. These are two different statements.

\subsection{Nitrogen pathway coverage differs between fuels}

Nitric oxide forms by several competing routes. The two mechanisms in use do not cover
the same set, and neither contains the other.

\begin{center}
\begin{tabular}{lcc}
\toprule
Route & \code{A2NOx_skeletal} (Jet-A) & \code{gri30} (LNG) \\
\midrule
Thermal (Zeldovich) & yes & yes \\
N$_2$O route & yes & yes \\
Prompt via NCN (modern) & \textbf{yes} & no \\
Prompt via HCN (older) & yes & yes \\
NNH route & \textbf{no} & \textbf{yes} \\
\bottomrule
\end{tabular}
\end{center}

Each mechanism happens to carry the route most important to its own regime: prompt NO
matters in the rich zones where Jet-A burns, and the NNH route matters in the lean
premixed conditions where LNG burns. That is fortunate, but it is not a reason to ignore
it.

\warningbox{Because the two fuels use mechanisms with different pathway coverage,
\emph{absolute} NO$_x$ values are not strictly comparable between Jet-A and LNG. The
design metric defined later compares each fuel only against itself, which avoids this
problem entirely.}

Rather than demanding a fixed species list, the code records which routes each mechanism
carries and attaches that record to every result. A hard requirement for NNH would
wrongly reject the Jet-A mechanism, which is otherwise the better of the two. Implemented
by \code{nox_pathway_coverage} in \code{src/fuelnozzle/crn/chemistry.py}.

\subsection{What is checked, and what merely warns}

A missing nitrogen submodel \emph{raises an error} rather than issuing a warning. Such a
mechanism reports exactly zero NO with no other symptom, so a warning could be missed and
the resulting zeros believed. The same applies to a fuel species absent from the
mechanism, and to an LNG component that the mechanism cannot represent: dropping a
component would silently change the fuel being modeled.

Conditions outside a mechanism's declared validity range produce warnings rather than
errors, because extrapolation is sometimes the only option and the user is entitled to
make that choice knowingly. GRI-Mech 3.0 was optimized against data up to roughly
10~atm; at combustor pressures of 10 to 40~atm every LNG result carries a
warning to that effect. Implemented by \code{validate_mechanism}.

\section{Air-Side Bookkeeping and Equivalence Ratio}

\subsection{Equivalence ratio}

The \emph{equivalence ratio} $\phi$ says how rich or lean a mixture is. It is the actual
fuel-to-air mass ratio divided by the ratio that would consume both exactly:

\begin{equation}
  \phi=\frac{\left(\dot m_{f}/\dot m_{a}\right)}
            {\left(\dot m_{f}/\dot m_{a}\right)_{\mathrm{stoich}}}
       =\frac{\dot m_{f}\,\mathrm{AFR}_{\mathrm{stoich}}}{\dot m_{a}}.
  \label{eq:equivalence-ratio}
\end{equation}

$\phi<1$ is lean, $\phi=1$ is stoichiometric, and $\phi>1$ is rich. Peak temperature, and
therefore most NO$_x$, occurs slightly rich of stoichiometric.

Computing $\phi$ by counting unburned fuel and air fails once reaction has started, since
the fuel is gone but the mixture has not changed identity. The code therefore uses an
\emph{element-based} definition, which counts carbon, hydrogen, and oxygen atoms rather
than molecules and so remains correct in burned, partially burned, and evaporating
mixtures alike. The related \emph{Bilger mixture fraction} $Z$ measures the same thing on
a scale running from 0 in pure air to 1 in pure fuel, and the two are related through the
stoichiometric mixture fraction $Z_{st}$:

\begin{equation}
  \phi=\frac{Z}{1-Z}\cdot\frac{1-Z_{st}}{Z_{st}},
  \qquad Z_{st}=\frac{1}{1+\mathrm{AFR}_{\mathrm{stoich}}}.
  \label{eq:phi-from-z}
\end{equation}

Both are evaluated by Cantera's own routines rather than reimplemented here. Verified
against hand calculation: methane gives $\mathrm{AFR}_{\mathrm{stoich}}=17.13$ and the
Jet-A surrogate C$_{11}$H$_{22}$ gives $14.69$, and Equation~\eqref{eq:phi-from-z}
reproduces a prescribed $\phi$ to six decimal places.

\subsection{Where the air goes}

Combustor air does not all enter at one place. It is divided among the dome, the primary
zone, quench jets, dilution jets, and liner cooling films:

\begin{equation}
  f_{\mathrm{dome}}+f_{\mathrm{primary}}+f_{\mathrm{quench}}
  +f_{\mathrm{dilution}}+f_{\mathrm{cooling}}=1.
  \label{eq:air-split}
\end{equation}

The code refuses a split that does not satisfy Equation~\eqref{eq:air-split} rather than
renormalizing it, because a split that does not sum to one usually means a station was
forgotten, and silently rescaling would hide the mistake.

\subsection{The inactive circuit, and why it is a design lever}

Jet-A and LNG have separate injectors. Air flows through \emph{both} injector air
passages at all times, including the one whose fuel is shut off. If the Jet-A passage
carries a share $a_J$ of the dome air and the LNG passage carries $a_L$, with
$a_J+a_L=f_{\mathrm{dome}}$, then the air setting the equivalence ratio near the flame
depends on which fuel is burning:

\begin{equation}
  f_{\mathrm{near}}(\text{fuel})=a_{\mathrm{active}}
  +\chi\,a_{\mathrm{idle}},
  \label{eq:near-field-air}
\end{equation}

where $\chi\in[0,1]$ is the fraction of the idle passage's unfueled air that actually
reaches the near field.

Equation~\eqref{eq:near-field-air} is the reason a single combustor may be able to run
rich for Jet-A and lean for LNG. Making $a_J$ small and $a_L$ large produces a rich Jet-A
near field and a lean LNG near field with no moving parts and no change in hole area.

\warningbox{The strength of this effect depends entirely on $\chi$, and a reactor network
cannot determine $\chi$ --- it is a mixing question that requires computational fluid
dynamics or an experiment. Every result that depends on it is reported with a sensitivity
across the full range $\chi\in[0,1]$, never at a single assumed value. At $\chi=0$ the
lever is at full strength; at $\chi=1$ it disappears entirely.}

Implemented by \code{AirSplit} and \code{resolve_air_streams} in
\code{src/fuelnozzle/crn/streams.py}.

\subsection{Reporting emissions}

Raw concentrations are not comparable between combustors, because a leaner combustor
dilutes its own exhaust. Two conventional corrections are applied. The \emph{emission
index} expresses output per unit of fuel burned,

\begin{equation}
  \mathrm{EI}_i=1000\,\frac{Y_i\,\dot m_{\mathrm{exhaust}}}{\dot m_{\mathrm{fuel}}}
  \quad[\mathrm{g/kg}],
  \label{eq:emission-index}
\end{equation}

and concentrations are quoted dry, with water removed, corrected to a reference oxygen
level of 15\%:

\begin{equation}
  \mathrm{ppmv}_{15\%}=10^{6}\,X_i^{\mathrm{dry}}\,
  \frac{20.9-15}{20.9-X_{\mathrm{O_2}}^{\mathrm{dry}}\times100}.
  \label{eq:o2-correction}
\end{equation}

Equation~\eqref{eq:o2-correction} is undefined when the measured oxygen approaches that
of air, which happens as combustion approaches extinction; the code raises rather than
returning a large meaningless number.

\section{Spray Boundary: Connecting the Nozzle to the Network}

A spray is described to the network by how much fuel arrives as vapor, and by a set of
\emph{droplet classes}: groups of droplets sharing a size, temperature, and velocity.
The number of droplets follows from the liquid mass flow and the mass of one droplet,

\begin{equation}
  \dot N_d=\frac{\dot m_{\ell}}{\rho_{\ell}\,\tfrac{4}{3}\pi r_0^{3}}.
  \label{eq:number-rate}
\end{equation}

For LNG, part of the fuel has already flashed to vapor inside the nozzle. Only the
remainder becomes droplets. Treating all of it as liquid would ask the model to perform
atomization work the flash has already done.

\subsection{Where the initial size comes from}

Three policies are available. The default takes the initial radius as half the Sauter
mean diameter from the nozzle model, which uses information the reference paper did not
have, since it solves the atomizer separately. The second follows the paper: start at the
orifice radius and let aerodynamic breakup find the size. The third is an explicit
user-supplied value.

\warningbox{The existing package refuses to report a Sauter mean diameter without a
hardware calibration. That refusal is inherited here. If no calibration exists, the
policy silently degrading to a fabricated size would be worse than useless, so instead it
falls back to the orifice-radius approach and emits a warning stating that
atomization-quality results must not be reported from it.}

\subsection{Representing a range of sizes}

A real spray is not one size. The largest droplets survive far longer than the mean and
are the ones that actually reach the flame, so representing only the mean understates how
much liquid arrives unevaporated. The Rosin-Rammler distribution describes the mass
fraction contained in droplets smaller than $D$ as

\begin{equation}
  Q(D)=1-\exp\!\left[-\left(\frac{D}{X}\right)^{n}\right],
  \label{eq:rosin-rammler}
\end{equation}

where $X$ is a characteristic diameter and $n$ measures how narrow the spray is. Its
Sauter mean diameter is

\begin{equation}
  D_{32}=\frac{X}{\Gamma\!\left(1-\tfrac{1}{n}\right)},
  \label{eq:rr-smd}
\end{equation}

which is inverted so that the classes reproduce the Sauter mean diameter requested.
Equation~\eqref{eq:rr-smd} diverges as $n\to1$, so a lower bound on $n$ is enforced
rather than allowing a distribution with no finite Sauter mean.

\section{Droplet Breakup}

A droplet moving quickly through gas is deformed by the pressure difference across it and
pulled back into shape by surface tension. The Taylor analogy treats that competition as a
mass on a spring with friction. Writing $y$ for the distortion scaled so that $y=1$ means
the deformation has reached the droplet radius,

\begin{equation}
  \frac{d^{2}y}{dt^{2}}
  =\frac{C_f}{C_b}\frac{\rho_g u_d^{2}}{\rho_{\ell}r_0^{2}}
  -\frac{C_k\sigma}{\rho_{\ell}r_0^{3}}y
  -\frac{C_d\mu_{\ell}}{\rho_{\ell}r_0^{2}}\frac{dy}{dt},
  \label{eq:tab}
\end{equation}

with $C_k=8$, $C_f=1/3$, $C_d=5$, $C_b=1/2$. The three terms are, in order, the
aerodynamic forcing, the surface-tension restoring force, and viscous damping. The
balance between forcing and surface tension is the Weber number,

\begin{equation}
  \mathrm{We}_g=\frac{\rho_g u_d^{2} r_0}{\sigma},
  \label{eq:weber}
\end{equation}

and the droplet breaks when the distortion reaches unity. Under steady forcing the
distortion oscillates about an equilibrium value $\mathrm{We}_c=\mathrm{We}_g/12$ with
amplitude $A$, so breakup is possible when $\mathrm{We}_c+A>1$.

\warningbox{John et al. print this criterion as $A+\mathrm{We}_g>1$. That cannot be what
they computed: for their own Fig.~3 case $\mathrm{We}_g\approx38$, so the criterion would
fire immediately and no breakup time would be resolvable, yet their figure shows a finite
$2\times10^{-5}$~s. The standard form $\mathrm{We}_c+A>1$ is used here and reproduces a
breakup time of the right order.}

The child radius follows from an energy balance between the parent's distortion energy
and the new surface created,

\begin{equation}
  r=\frac{r_0}{1+\tfrac{8K}{20}y^{2}
  +\dfrac{\rho_{\ell}r_0^{3}\dot y^{2}}{\sigma}\dfrac{6K-5}{120}},
  \qquad K=\tfrac{10}{3}.
  \label{eq:child-radius}
\end{equation}

One pass rarely reaches a stable size, so each child becomes the new parent and the test
repeats until the Weber number is too small to sustain further breakup.

\warningbox{The Taylor analogy describes a droplet torn apart from outside by aerodynamic
forces. A flashing LNG droplet is burst from within by vapor generation, which is a
different mechanism entirely. Breakup is therefore not applied in the flashing regimes;
the size comes from the Tier 3 flash-spray result instead, and the decision is recorded
as a warning rather than left implicit.}

\section{Droplet Evaporation}

Evaporation can be limited by either of two different things, and which one applies
decides the equations.

\subsection{Diffusion limited}

Below the boiling point, fuel vapor must diffuse away from the surface before more can
leave. The surface vapor mass fraction follows from the vapor pressure,

\begin{equation}
  Y_s=\frac{W_F}{W_F+W_{\mathrm{mix}}\left(\dfrac{p_g}{p_v(T_d)}-1\right)},
  \label{eq:surface-fraction}
\end{equation}

and the difference between surface and far-field concentration drives the transfer,

\begin{equation}
  B_d=\frac{Y_s-Y_\infty}{1-Y_s},
  \qquad
  \mathrm{Sh}=\left(2+0.6\,\mathrm{Re}_d^{1/2}\mathrm{Sc}^{1/3}\right)
  \frac{\ln(1+B_d)}{B_d},
  \label{eq:spalding}
\end{equation}

\begin{equation}
  \frac{dr}{dt}=-\frac{\rho_g D}{2\rho_{\ell}r}B_d\,\mathrm{Sh}.
  \label{eq:evaporation}
\end{equation}

In still gas this reduces to $d(r^{2})/dt=\mathrm{constant}$, the classical
\emph{$D^2$ law}: the square of the diameter falls linearly with time. Any evaporation
model must reproduce it, and this one is verified to do so --- but only after the droplet
has finished warming up, since the film temperature moves during that transient.

\subsection{Boiling limited}

Equation~\eqref{eq:surface-fraction} gives $Y_s=1$ when the vapor pressure reaches the
ambient pressure, at which point $B_d$ becomes infinite and the diffusion description
fails. That is the boiling point, and it is exactly where the second branch takes over.

A boiling droplet cannot get any hotter: its temperature is pinned at saturation and
every joule that arrives goes into the phase change,

\begin{equation}
  \frac{dm_d}{dt}=-\frac{\dot Q_{\mathrm{conv}}}{L_{\mathrm{vap}}}.
  \label{eq:boiling}
\end{equation}

\warningbox{This branch is the single most important addition to the reference model.
Flashing LNG arrives above its saturation temperature, so it is \emph{never} diffusion
limited. Applying Equation~\eqref{eq:evaporation} to it would be wrong at the first step.}

\section{Droplet Heating and Energy Accounting}

Heat reaches the droplet by convection,

\begin{equation}
  \dot Q_{\mathrm{conv}}
  =\frac{A_d\,\beta_{\mathrm{spray}}\,\mathrm{Nu}_d\,k(T_g-T_d)}{2r},
  \qquad
  \mathrm{Nu}_d=\left(2+0.6\,\mathrm{Re}_d^{1/2}\mathrm{Pr}_d^{1/3}\right)
  \frac{\ln(1+B)}{B},
  \label{eq:convective-heat}
\end{equation}

with gas properties evaluated at the film temperature
$\bar T=(T_g+2T_d)/3$, weighted toward the droplet because most of the resistance sits in
the thin layer against the surface.

Part of that heat drives the phase change and the rest warms the liquid:

\begin{equation}
  c_{\ell}m_d\frac{dT_d}{dt}
  =\dot Q_{\mathrm{conv}}+\frac{dm_d}{dt}L_{\mathrm{vap}}.
  \label{eq:droplet-heating}
\end{equation}

Since $dm_d/dt$ is negative, the second term subtracts, as it must.

\subsection{The bookkeeping that is easy to get wrong}

Three statements must hold together, and stating them separately is what keeps the energy
balance honest:

\begin{enumerate}[label=(\roman*)]
  \item the gas \emph{loses} $\dot Q_{\mathrm{conv}}$;
  \item the droplet \emph{gains} $\dot Q_{\mathrm{conv}}$, spends
        $|dm_d/dt|L_{\mathrm{vap}}$ on vaporization, and warms with the remainder;
  \item the vapor enters the gas carrying the enthalpy it has \emph{at the droplet
        temperature} $T_d$, not at the gas temperature.
\end{enumerate}

If item (iii) is neglected and the vapor is added at gas temperature, the latent heat is
counted twice and the reactor runs hot. Equation~\eqref{eq:droplet-heating} is verified
directly: the convective heat equals the sum of the latent and sensible terms to ten
significant figures.

\section{Reactor Types and the Cantera Implementation}

A \emph{reactor} here is a region of the combustor treated as internally uniform. Which
idealization it is given states what is being assumed about mixing inside it. A
\emph{perfectly stirred reactor} assumes anything entering is instantly spread
throughout, which suits recirculation and flame zones where turbulence is violent. A
\emph{plug flow reactor} assumes the flow marches along without mixing backwards, which
suits the post-flame length where axial velocity dominates. An \emph{evaporator} and a
\emph{mixer} are stirred reactors that additionally contain droplets.

Plug flow is represented as a chain of stirred reactors rather than by Cantera's
\code{FlowReactor}, so that it participates in the same simultaneous network solution as
everything else. The number of links is a numerical parameter whose convergence must be
demonstrated, not assumed.

The residence time of a zone is

\begin{equation}
  \tau_i=\frac{\bar\rho_i V_i}{\dot m_i}.
  \label{eq:residence-time}
\end{equation}

\warningbox{Because a reactor stands for a coarse region with real internal variation,
Equation~\eqref{eq:residence-time} gives an \emph{effective zone-scale} residence time.
It is not the transit time of any particular parcel of fluid, and should not be read as
one.}

\section{Network Assembly, Mass Balance, and Steady Solution}

\subsection{The splits will not balance, and what to do about it}

Flow splits supplied by a user, or extracted from a computational fluid dynamics
solution, will not conserve mass exactly. John et al. report imbalances up to 9.3\% on a
single reactor. For each zone, what enters must equal what leaves:

\begin{equation}
  \dot m_{i,\mathrm{in}}+\sum_j \dot m_{ji}
  =\dot m_{i,\mathrm{out}}+\sum_j \dot m_{ij}.
  \label{eq:reactor-mass-balance}
\end{equation}

Collecting the reactor-to-reactor flows into a vector $\mathbf z$, that is a linear
system $\mathbf{Az}=\mathbf b$, which the supplied flows miss by a residual
$\mathbf r=\mathbf b-\mathbf{Az}$. We want the \emph{smallest} change $\Delta\mathbf z$
that repairs it:

\begin{equation}
  \Delta\mathbf z^{*}=\arg\min\|\Delta\mathbf z\|_2^{2}
  \quad\text{subject to}\quad \mathbf A\,\Delta\mathbf z=\mathbf r.
  \label{eq:min-norm-correction}
\end{equation}

There are normally more flows than reactors, so the system is underdetermined and
Equation~\eqref{eq:min-norm-correction} is exactly what a least-squares solve returns.
Choosing the smallest change rather than rescaling means no single flow absorbs the
entire error, and flows that were already consistent are barely touched.

The rows of $\mathbf A$ sum to zero, since every internal flow leaves one reactor and
enters another. A solution therefore exists only if the external streams balance
globally. When they do not, the boundary conditions themselves are wrong and no
redistribution can help, so the code raises rather than returning a plausible-looking
repair.

\subsection{Solving the network}

Recirculation is what makes this a graph rather than a chain: hot products flowing back
toward the injector are what hold the flame, so a zone's inlet can depend on its own
downstream. All zones are therefore handed to a single Cantera \code{ReactorNet} and
solved together, which handles the recycle without any iteration on our part.

The steady state is reached by \emph{marching in time} until nothing moves any more,
rather than by calling Cantera's steady-state solver.

\warningbox{This reverses the original plan, for a measured reason. On a four-reactor,
71-species recirculating network, \code{advance_to_steady_state} ran for 210~s and then
failed. The same case marches to a converged answer, with an identical peak temperature
of 1881~K, in 0.86~s. Marching is also bounded and observable: it integrates in chunks of
two residence times, stops when every reactor temperature moves less than 0.05~K, and
gives up after two hundred residence times with a warning. The cost is that element
conservation closes to about $10^{-7}$ rather than $10^{-11}$, reflecting that tolerance.}

\subsection{Detecting a solution that is not a combustor}

A network initialized cold converges to a perfectly valid steady state in which nothing
is burning. It solves the equations and it is useless. Reactors are therefore started hot
and equilibrated, and any converged state whose peak temperature is within 50~K of the
inlet air is reported as extinguished, because emissions computed from it are
meaningless.

\section{Coupling the Droplets to the Network}

Evaporation depends on gas temperature, and gas temperature depends on evaporation.
Solving both at once would place a stiff droplet system inside a stiff chemistry system.
Instead the two are alternated: the droplets are marched with the gas held fixed, the
network is solved with the resulting sources held fixed, and the process repeats with
only part of each update applied.

The partial step matters. Evaporation and temperature reinforce each other, since hotter
gas evaporates more fuel, which burns and makes the gas hotter still, so taking the full
step invites oscillation between a flooded and a starved state.

Two details keep this honest. The vapor enters each zone at the \emph{droplet}
temperature, and the heat the droplets absorbed is removed from the gas; without the
second, injecting vapor would hand the gas back the latent heat it had just spent. And
the network is rebuilt rather than edited on each pass, because evaporating fuel changes
the total gas mass flow and therefore the mass balance itself.

Droplets surviving the whole spray path are reported as \emph{liquid carryover}. This is
a design failure rather than a rounding detail: fuel that never evaporates never burns.

\section{Combustor Architecture Templates}

The architectures differ in one decision, how much air meets the fuel before it burns,
and that decision is what the dual-fuel study exists to make.

\subsection{Rich-quench-lean}

Burn the fuel rich, where it is too fuel-heavy to get very hot, then add air fast enough
to cross stoichiometric quickly and finish lean.

\warningbox{The quench is modeled as a chain of stages, not a single mixing point. This
is not cosmetic. Nitric oxide is produced \emph{while} the mixture crosses stoichiometric
on its way from rich to lean, and a single perfectly mixed point would jump straight over
that crossing and hide the dominant source entirely.}

Quench air is distributed along the stages by a declared schedule. Front-loading shortens
the traverse but risks quenching the reaction; rear-loading does the opposite. The
sensitivity to stage count and schedule is a required output, not an option.

\subsection{Lean premixed and lean direct injection}

Lean premixed mixes fuel and air completely before burning, so no part of the flame is
hot enough to make much nitric oxide. Lean direct injection sprays into a lean dome with
no premixing passage at all.

\warningbox{Chemistry is \emph{not} frozen in the premixer, even though the whole idea of
a premixer is that nothing reacts inside it. Freezing it would guarantee the assumption
instead of testing it. Left reacting, the premixer ignites in the model precisely when
residence time exceeds ignition delay, which is the real failure mode of the
architecture. Lean direct injection has no premixing passage and therefore nowhere for
this to happen, which is exactly why it is the natural comparison for Jet-A at landing
and take-off.}

\subsection{The idle passage as a design lever}

Because the two fuel circuits are separate hardware, air flows through \emph{both}
injector passages at all times, including the one whose fuel is shut off. If the Jet-A
passage carries a share $a_J$ of the dome air and the LNG passage $a_L$, the air setting
the equivalence ratio near the flame is

\begin{equation}
  f_{\mathrm{near}}=a_{\mathrm{active}}+\chi\,a_{\mathrm{idle}},
  \label{eq:near-field-template}
\end{equation}

where $\chi$ is the fraction of idle-passage air reaching the near field. Every template
routes this explicitly. Making $a_J$ small and $a_L$ large gives a rich Jet-A near field
and a lean LNG near field with no moving parts.

\warningbox{At $\chi=0$ the lever is at full strength; at $\chi=1$ it disappears
entirely. A reactor network cannot determine $\chi$, so every result depending on it is
reported across its full range rather than at one assumed value.}

\subsection{Merging mechanisms}

Combining a fuel mechanism with a nitrogen submodel is sometimes necessary, and the tool
checks such a combination without performing it, because the choice changes the answer
and belongs to whoever is accountable for it. Checked against the two mechanisms actually
in use, it finds six species with inconsistent thermodynamic data, 223 duplicate
reactions, and a species spelled \code{CH2*} in one file and \code{CH2(S)} in the other.
The last is an error rather than a warning: a merged mechanism would carry two separate
pools of what is physically one species.

\section{Autoignition and Flashback Screening}

A premixing passage exists so that fuel and air are uniform before they burn. It works
only if the mixture takes longer to ignite on its own than it takes to cross the passage.
The margin is a ratio of two times,

\begin{equation}
  M=\frac{\tau_{\mathrm{ign}}}{\tau_{\mathrm{res}}},
  \label{eq:autoignition-margin}
\end{equation}

and a design guidance floor of $M\geq4$ is used as an editable input rather than a hidden
constant. When $M<1$ the mixture ignites inside the injector, and the result is reported
as an error rather than a number.

\subsection{Why cold fuel buys margin}

Fuel that has flashed arrives cold, and whatever is still liquid draws its latent heat out
of the air. Both effects lower the premixed temperature, and because ignition delay
depends on temperature exponentially, a modest cooling buys a large increase in margin.
Measured here with air at 800~K and 20~atm:

\begin{center}
\begin{tabular}{lccc}
\toprule
Case & $T_{\mathrm{mix}}$ & Drop & of which latent \\
\midrule
Jet-A vapour at 470~K & 774.1~K & 25.9~K & 0.0~K \\
Jet-A, 30\% liquid & 771.5~K & 28.5~K & 2.6~K \\
LNG vapour at 150~K & 747.2~K & 52.8~K & 0.0~K \\
LNG, 30\% liquid & 742.8~K & 57.2~K & 4.4~K \\
\bottomrule
\end{tabular}
\end{center}

LNG cools the mixture about twice as much as Jet-A. Note that the latent contribution is
only about 4~K of a 57~K drop: the benefit is overwhelmingly the sensible enthalpy of cold
vapour, not the phase change. It would be easy, and wrong, to assume the opposite.

\subsection{What that is worth}

\begin{center}
\begin{tabular}{lcccc}
\toprule
Fuel & $T_{\mathrm{mix}}$ & $\tau_{\mathrm{ign}}$ & $M$ at 1~ms & $M$ at 3~ms \\
\midrule
Jet-A & 774~K & 4.04~ms & 4.0 (safe) & 1.3 (marginal) \\
LNG & 747~K & no ignition in 10~s & unbounded & unbounded \\
\bottomrule
\end{tabular}
\end{center}

Jet-A tolerates roughly one millisecond of premixing at landing and take-off conditions.
LNG is effectively inert on premixer timescales. This is the quantitative basis for
considering a lean premixed LNG cruise alongside a lean direct injection or
rich-quench-lean Jet-A landing and take-off.

\warningbox{The flashback screen is \emph{suppressed} unless a laminar flame speed is
supplied. Obtaining one requires a flame calculation this tool does not perform, and
inventing a value would place an unmarked guess underneath a safety screen. The screen
reports nothing rather than something unfounded.}

\section{Cryogenic Thermal Fuel Management}

LNG leaves the tank near 110~K and must reach the injector warm enough to flash usefully.
Fuel temperature is therefore a design variable, pulled in four directions at once:
enough superheat at the nozzle for the intended flash regime; enough subcooling upstream
that the feed line does not boil; enough autoignition margin; and within the heat actually
available to dump into the fuel.

\subsection{The window, and what sets its width}

Superheat at the injector and subcooling in the line compete for the same resource: the
temperature gap between saturation at feed pressure and saturation at chamber pressure.
With a 20~bar chamber:

\begin{center}
\begin{tabular}{lccc}
\toprule
Feed pressure & Saturation gap & Feasible $T_{\mathrm{fuel}}$ & Width \\
\midrule
2.5~MPa & 6.1~K & \textbf{none} & --- \\
3.0~MPa & 11.4~K & 171.0--172.0~K & 1.0~K \\
4.0~MPa & 20.2~K & 171.0--181.0~K & 10.0~K \\
\bottomrule
\end{tabular}
\end{center}

\warningbox{Below roughly 3~MPa feed pressure there is \emph{no} feasible fuel temperature
at all. No amount of heating fixes this, because the gap itself is too narrow. Pump
pressure is what buys thermal design freedom, which connects this window directly to the
pressure budget the package already computes.}

Above the critical pressure of the mixture there is no saturation at all. Supercritical
injection is a legitimate design choice, but the subcooling and superheat constraints are
defined only below the critical point and do not apply to it; the code says so explicitly
rather than failing obscurely.

\subsection{The circuit that is not running}

A dual-fuel nozzle always has one circuit hot and stagnant. With no flow there is no
convective cooling, so the idle fuel sits near wall temperature. Jet-A left in a hot
passage forms deposits that block the injector; LNG left in one boils, and restart then
delivers vapour rather than liquid. Either can veto an otherwise attractive packaging, and
it is cheaper to discover here than on a test stand.

\section{Verification Coverage}

Verification asks whether the software solves the equations it claims to. It is kept
separate from validation, which asks whether those equations describe reality.

\subsection{Limits and identities}

The residence-time definition reproduces the identity the reference paper states.
A stirred reactor given a long residence time approaches the equilibrium state. Shrinking
it below a critical size extinguishes it. With no fuel, air passes through unchanged.
Droplet evaporation recovers the $D^2$ law once the warm-up transient is excluded, and the
boiling branch reproduces $\dot m=\dot Q/L$ exactly. The droplet energy split closes to ten
significant figures. Element balances close to about $10^{-7}$.

\subsection{Convergence}

A plug flow zone represented as a chain of stirred reactors converges with the number of
links. The droplet integration is insensitive to solver tolerance across four orders of
magnitude. The ignition-delay table changes by under 15\% when its temperature grid is
refined fourfold.

\subsection{Comparison with the reference paper}

The purpose is directional: to confirm the implementation reproduces the paper's trends
and conclusions, not its absolute numbers. Two of its quantitative anchors are only
partly reproducible, and this is stated rather than smoothed over.

\emph{Breakup and evaporation timescales.} Breakup agrees to an order of magnitude, at
$3.4\times10^{-5}$~s against their $2\times10^{-5}$~s. Evaporation does not: we obtain
0.80~s against their reported 6~ms at 350~K. Their case is underspecified, giving neither
ambient pressure nor liquid properties, and 6~ms is difficult to reconcile with realistic
kerosene volatility. What is reproduced is their \emph{conclusion}: breakup is faster than
evaporation by more than two orders of magnitude at 350~K, the separation narrows with
temperature, and it still exceeds one order of magnitude at 900~K. That conclusion is what
justifies treating the two sequentially, which is what the implementation depends on.

\emph{Liquid against gaseous networks.} This is reproduced cleanly. The same topology and
flows, run once with droplets and once with the fuel gaseous and perfectly mixed:

\begin{center}
\begin{tabular}{lcccc}
\toprule
 & $\phi$ range & Peak $T$ & Exit $T$ & NO \\
\midrule
With spray & 0.514--1.714 & 2453~K & 1926.6~K & 1868~ppm \\
Gaseous, premixed & 0.514--0.514 & 1928~K & 1928.0~K & 57~ppm \\
\bottomrule
\end{tabular}
\end{center}

Exit temperature differs by 0.07\%. Nitric oxide differs by a factor of 33. That asymmetry
is the paper's entire argument: both methods get the global energy balance right, and only
one gets the emissions right, because nitric oxide responds exponentially to the hot
pockets that spray heterogeneity creates. Their reported contrast, a gaseous network
underpredicting by 54--91\% with peak temperatures near 1800~K against above 2200~K with
spray, matches this in both direction and character.

\warningbox{Achieving this required the gaseous comparison to be genuinely premixed, with
fuel distributed to every zone in proportion to its air. Injecting the vapour only at the
dome retains the air-staging heterogeneity and is not the paper's comparison; the two cases
then differ by a few percent instead of a factor of thirty.}

\subsection{What is not claimed}

Absolute nitric oxide is not claimed to match the paper. Their reactor volumes,
connectivity, spray path lengths, and propagation fractions are unpublished, so their
calibrated network cannot be rebuilt. Their chemistry also differs: they used a mechanism
whose oxidation base carries no nitrogen chemistry, and the paper does not state which
nitrogen submodel they appended.

Absolute emission levels from this implementation are also not validated, and the reason
is now understood rather than merely suspected.

\subsection{How fast the quench is decides the answer}

The rich-quench-lean optimum initially fell at $\phi_{\mathrm{rich}}=1.22$, richer than
classical practice indicates, on a curve so flat it barely had a minimum. A sensitivity
study established that the cause was an unrealistically long quench, not the reactor model.

Discretization first. Splitting the quench into stages matters, and the size of the effect
is worth stating plainly: a \emph{single} perfectly mixed quench point underpredicts
nitric oxide by more than half, because it jumps straight over the stoichiometric crossing
where most of it is made. The result converges to about 1\% by twelve stages.

Quench \emph{speed} matters more:

\begin{center}
\begin{tabular}{lcc}
\toprule
Quench residence time & Optimum & Minimum $\mathrm{EI}_{\mathrm{NOx}}$ \\
\midrule
$\sim$0.9~ms (realistic) & $\phi_{\mathrm{rich}}=1.35$, sharp & 68.5 \\
$\sim$4.6~ms & $\phi_{\mathrm{rich}}=1.22$, flat & 130.6 \\
$\sim$18~ms & no interior minimum & 176.9 \\
\bottomrule
\end{tabular}
\end{center}

Real quench jets penetrate and mix in roughly half a millisecond to a millisecond, and
that speed is the entire point of the architecture. At a realistic quench the optimum
lands at 1.35 with a pronounced minimum, meeting the sanity anchor. A slow quench dwells
near stoichiometric, inflates nitric oxide, and progressively destroys the very optimum
the architecture exists to exploit.

\warningbox{The optimum itself exists only when there is enough quench \emph{air}, not
merely enough quench speed. Measured at four quench-air fractions, an interior optimum is
present at 0.30 and 0.38 of the total air but absent at 0.15 and 0.20, where the curve
inverts and $\phi_{\mathrm{rich}}\approx1.4$ becomes the worst choice rather than the
best. Rich-quench-lean depends on crossing stoichiometric quickly, and the crossing needs
both time and air. Any statement about a rich-zone optimum must name the quench-air
fraction it assumes.}

\warningbox{Absolute emission indices remain roughly twice what real rich-quench-lean
hardware achieves. The zones here are adiabatic, and a perfectly stirred quench still
mixes more slowly than real jets even at one millisecond. The optimum's \emph{location}
and the \emph{trends} are supported; the absolute \emph{level} is not.}

\section{Known Deviations and Technical Debt}

The numbered deviations table lives in \code{docs/CRN_IMPLEMENTATION_LOG.md} and is the
authoritative record. The ones that change how a result should be read are:

\begin{itemize}
  \item Jet-A uses two mechanisms, one for the network and one for ignition delay,
        because a high-temperature-only mechanism overpredicts ignition delay by 510$\times$
        at 700~K and would declare an unsafe premixer safe.
  \item The premixer is modelled as reacting rather than chemically frozen, so that its
        real failure mode is visible in the solution.
  \item The breakup criterion follows the standard Taylor analogy rather than the form
        printed in the paper, which cannot be what they computed.
  \item Steady state is reached by marching in time rather than by Cantera's steady
        solver, which failed on a recirculating network. Element conservation closes to
        about $10^{-7}$ rather than $10^{-11}$ as a result.
  \item The flashback screen and all droplet-size outputs are suppressed rather than
        estimated when the inputs they need are absent.
\end{itemize}

\appendix

\section{Cantera Interface Used by This Work}

Pinned to \textbf{Cantera 3.2.0}. Verified present in the installed version rather than
taken from documentation, because the reactor and flow-device interface changed between
Cantera 2.6 and 3.x.

\begin{center}
\begin{tabular}{lll}
\toprule
Purpose & Cantera name & Note \\
\midrule
Gas mixture and kinetics & \code{Solution} & One per reactor; never shared \\
Fixed boundary condition & \code{Reservoir} & Air inlets, fuel vapor, exhaust \\
Perfectly stirred reactor & \code{IdealGasConstPressureMoleReactor} & Mole-based; see below \\
Alternative reactor forms & \code{IdealGasReactor}, \code{IdealGasMoleReactor} & \\
Custom governing equations & \code{ExtensibleIdealGasConstPressureReactor} & For coupled droplets \\
Steady plug flow & \code{FlowReactor} & Not \code{PlugFlowReactor}; see below \\
Prescribed stream & \code{MassFlowController} & Reactor-to-reactor flows \\
Pressure closure & \code{PressureController} & Network exit \\
Pressure-driven flow & \code{Valve} & \\
Heat transfer & \code{Wall} & Liner heat loss \\
Surface chemistry & \code{ReactorSurface} & Not currently used \\
Network assembly and solution & \code{ReactorNet} & \\
Steady solution & \code{ReactorNet.advance_to_steady_state} & \\
Time integration & \code{ReactorNet.advance}, \code{.step} & \\
Large-mechanism speedup & \code{AdaptivePreconditioner} & Mole reactors only \\
Stoichiometry & \code{Solution.stoich_air_fuel_ratio} & \\
Equivalence ratio & \code{Solution.equivalence_ratio} & Element-based \\
Mixture fraction & \code{Solution.mixture_fraction} & Bilger \\
Profile storage & \code{SolutionArray} & \\
\bottomrule
\end{tabular}
\end{center}

Two findings from that verification are worth recording.

\emph{Mole-based reactors are not a stylistic preference.} Sparse Jacobian
preconditioning, which is what makes a network of 70-species reactors solve in seconds
rather than minutes, is supported only for \code{IdealGasMoleReactor} and
\code{IdealGasConstPressureMoleReactor}. That determines the default reactor type.

\emph{There is no \code{PlugFlowReactor} class.} The Cantera documentation describes a
``Plug Flow Reactor'', but the class is named \code{FlowReactor}. A plug flow zone is
represented here by a series of stirred reactors instead, which participates in the
simultaneous network solution; \code{FlowReactor} is reserved as an independent
cross-check.

\section{Building This Report}

From the repository root:

\begin{verbatim}
pixi run doc-crn
\end{verbatim}

which runs Tectonic on this file. \code{pixi run doc} builds both this report and the
nozzle technical reference.

\end{document}
